\subsection{Proofs for Section~\ref{sec:canonical}}\label{dfr:ssec:canonical}

The notation is that of Section~\ref{sec:canonical}.

\begin{proof}[Proof of Theorem~\ref{can:thm:monomial}]
Since $k!=0$ in $K$, $p=\operatorname{char}K$ satisfies $2\le p\le k$. Let $\nu\in S_k$ be a cycle of length $p$ and $g$
the permutation matrix with $ge_i=e_{\nu(i)}$, which is monomial. Corollary~\ref{can:cor:p-cycle}(2) excludes (1)--(3)
already for this $g$. For the algebraic powers, let $\mathsf H$ be the group generated by $\nu$, acting on $I=[k]$, so
that $V=K^k$. As in the proof of Corollary~\ref{can:cor:p-cycle}(2), a map equivariant for $g$ is
$\mathsf H$\nobreakdash-equivariant, and $\bar{\mathsf H}_{[k]}=\bar{\mathsf H}_{\mathbf m}$ has order $p$ for the
$k$\nobreakdash-multiset $\mathbf m$ with value $1$ on $[k]$; so Theorem~\ref{can:thm:permutation}(1)--(3) excludes such
maps.
\end{proof}

\begin{proof}[Proof of Corollary~\ref{can:cor:p-cycle}]
(1) We treat (1) and (2) of Theorem~\ref{can:thm:permutation}; for (3) replace $T$ by $\operatorname{supp}\mathbf m$ and
the stabilizers of $T$ by those of $\mathbf m$. The first statement implies the second by restriction. If $h\in\mathsf H$
maps a $k$\nobreakdash-subset $T$ onto itself and induces on it a permutation of order $p=\operatorname{char}K$, then $p$
divides $\abs{\bar C_T}$ for the group $C$ generated by $h$, and Theorem~\ref{can:thm:permutation} excludes a
$C$\nobreakdash-equivariant map; so the second statement implies the third. If the first fails,
Theorem~\ref{can:thm:permutation} gives a $k$\nobreakdash-subset $T$ such that $p=\operatorname{char}K>0$ divides
$\abs{\bar{\mathsf H}_T}$, and by Cauchy's theorem some element of $\mathsf H_T$ induces on $T$ a permutation of order
$p$; so the third fails. If $\operatorname{char}K=0$, all $\abs{\bar{\mathsf H}_T}$ and $\abs{\bar{\mathsf H}_{\mathbf m}}$
are nonzero in $K$.

(2) Let $\mathsf H$ be the group generated by $\nu$, acting on $I=[k]$ and on $E_0$ by $he_i=e_{h(i)}$; these operators
are bounded since the $e^*_i$ are, and satisfy~\eqref{can:eq:compatible} since
$e^*_i(he_j)=\delta_{i,h(j)}=e^*_{h^{-1}(i)}(e_j)$. For $T=I$ and the $k$\nobreakdash-multiset $\mathbf m$ with value $1$
on $I$, $\mathsf H_T=\mathsf H_{\mathbf m}=\mathsf H$, and $\bar{\mathsf H}_T=\bar{\mathsf H}_{\mathbf m}$ has order $p$,
which is $0$ in $K$. An equivariance identity for $g$ holds for its powers, since $(g^j)^{\otimes k}=(g^{\otimes k})^j$
and $\Lambda^k_\pi(g^j)=(\Lambda^k_\pi g)^j$, hence for all of $\mathsf H$, which Theorem~\ref{can:thm:permutation}(4)
excludes.
\end{proof}

\begin{proof}[Proof of Proposition~\ref{can:prop:asymmetry}]
(1) This is Proposition~\ref{prop:diag}(1) applied to $Q$. With $u_S:=\sum_{i\in S}u_i$, every $k$\nobreakdash-linear
$\Phi$ satisfies $\sum_{S\subseteq[k]}(-1)^{k-\abs S}\Phi(u_S,\dots,u_S)=\sum_\sigma\Phi(u_{\sigma(1)},\dots,u_{\sigma(k)})$,
as in~\eqref{can:eq:polarization}. Apply this to $Q(\cdot)\circ\omega_E$ and to
$(u_1,\dots,u_k)\mapsto\bigl((x_b)\mapsto u_1x_1\wedge\dots\wedge u_kx_k\bigr)$, which agree on the diagonal since
$Q(u,\dots,u)=\Lambda^k_\pi u$: the right sides give
$\sum_\sigma Q(u_{\sigma(1)},\dots,u_{\sigma(k)})\circ\omega_E=\Lsym(u_1,\dots,u_k)\circ\omega_E$, and pure wedges span
$\Lambda^k_\pi E$.
(2) A symmetric $Q$ would give $\Lsym=k!\,Q=0$ by~(1). Let $f_1,\dots,f_k$ be the functionals, choose
$x_1,\dots,x_k\in E$ with $f_i(x_j)=\delta_{ij}$ and, since $\Lambda^k_\pi D$ is dense in $\widehat\Lambda^k_\pi D\ne0$,
$d_1,\dots,d_k\in D$ with $d_1\wedge\dots\wedge d_k\ne0$. For $u_i=f_i(\cdot)\,d_i$ only $\sigma=\id$ contributes
to~\eqref{can:eq:sym}, so $\Lsym(u_1,\dots,u_k)(x_1\wedge\dots\wedge x_k)=d_1\wedge\dots\wedge d_k\ne0$.
(3) By~\eqref{can:eq:signs},
$\wedge^k_D=\rho\circ\operatorname{alt}^k_D\circ\wedge^k_D=\sum_\sigma\sgn(\sigma)\rho\circ\tau_\sigma=k!\,\rho=0$, so
$\Lambda^k_\pi D=\wedge^k_D(\otimes^k_\pi D)=0$.
(4) For $w\in\Lambda^k_\pi E$, \eqref{can:eq:signs} gives
$\operatorname{alt}^k_E(w)=\operatorname{alt}^k_E\widehat\wedge^k_Es(w)=\sum_\sigma\sgn(\sigma)\tau_\sigma s(w)=k!\,s(w)=0$.
With $f_i$ and $x_j$ as in~(2), the bounded functional on $\widehat\otimes^k_\pi E$ linearizing
$(y_a)\mapsto\prod_af_a(y_a)$ takes the value $\sum_\sigma\sgn(\sigma)\prod_af_a(x_{\sigma(a)})=1$ at
$\operatorname{alt}^k_E(x_1\wedge\dots\wedge x_k)$, a contradiction.
\end{proof}
